% Generated draft: requires mathematical and source review.
\documentclass[11pt]{article}
\usepackage[margin=25mm]{geometry}
\usepackage{fontspec}
\setmainfont{Latin Modern Roman}
\usepackage{amsmath,amssymb,mathtools,unicode-math}
\setmathfont{Latin Modern Math}
\usepackage{enumitem,needspace,xurl,hyperref}
\setlength{\parindent}{0pt}
\setlength{\parskip}{5pt}
\setlength{\emergencystretch}{3em}
\newcommand{\N}{\mathbb{N}}
\newcommand{\Z}{\mathbb{Z}}
\newcommand{\Q}{\mathbb{Q}}
\newcommand{\R}{\mathbb{R}}
\newcommand{\C}{\mathbb{C}}
\newcommand{\F}{\mathbb{F}}
\newcommand{\A}{\mathbb{A}}
\newcommand{\PP}{\mathbb{P}}
\newcommand{\E}{\mathbb{E}}
\newcommand{\eps}{\varepsilon}
\newcommand{\dd}{\,\mathrm{d}}
\begin{document}
\section*{UnsolvedMath ID: 30000535}
\textbf{Source problem code:} catalogue entry
\subsection{Definitions and mathematical statement}
Let \(\mathbb K\in\{\R,\C\}\) and \(G=SL_3(\mathbb K)\). A lattice \(\Gamma\leq G\) is a discrete subgroup of finite Haar covolume. It is noncocompact when \(G/\Gamma\) is noncompact. A matrix \(u\in G\) is unipotent when \((u-I_3)^3=0\); a unipotent subgroup in the lattice means a subgroup all of whose elements are unipotent matrices, not a positive-dimensional continuous subgroup contained in a discrete lattice.

The conjecture asks whether every such noncocompact \(\Gamma\) has a subgroup \(\Gamma_0\leq\Gamma\) of finite index and finitely many unipotent subgroups \(U_1,\ldots,U_m\leq\Gamma_0\) with
\[
 \Gamma_0=U_1U_2\cdots U_m.
\]
Explicitly, the quantifier order is
\[
 \forall\Gamma\ \exists\Gamma_0\ \exists m\geq1\ \exists U_1,\ldots,U_m
 \ \forall\gamma\in\Gamma_0\ \exists(u_1,\ldots,u_m)\in\prod_iU_i,
 \quad\gamma=u_1\cdots u_m.
\]
The ordered list of subgroups is fixed for \(\Gamma_0\), and the number of factors is independent of \(\gamma\). It may depend on the lattice. Repeated subgroups in the list are allowed; no uniqueness of the factors, normality or commutation of the \(U_i\), cyclicity, or bound uniform over all lattices is asserted. Finite index does not require \(\Gamma_0\) normal.

This bounded-product statement is stronger than saying that unipotent subgroups generate \(\Gamma_0\) with arbitrarily long words. It also does not require the factors to be entire real or complex unipotent algebraic groups; if such an ambient description is intended, the actual factors in the product must be the appropriate subgroups of the lattice. The source gives no specific root-subgroup list or arithmetic parametrization, and none is inserted.

The literature assessment checked 2026-08-21 reports bounded-generation results for broad families and reductions of remaining cases to the \(SL_3\) unipotent-generation conjecture, while saying the full statement is not settled. It does not specify those families or reduction hypotheses. This partial-progress report, the partially-solved label and editorial verification note remain dated metadata, not independently verified current results.
\subsection{Short English statement}
Must every noncompact finite-covolume lattice in \(SL_3(\R)\) or \(SL_3(\C)\) become a fixed finite product of unipotent subgroups after passing to finite index? The fixed factor bound is per lattice and independent of the element, so ordinary unipotent generation is insufficient.
\subsection{Sources}
\begin{itemize}
\item Local UnsolvedMath registry, permanent ID 30000535, source code catalogue entry, ``Bounded Unipotent Generation of Noncocompact \(SL_3\) Lattices,'' including original conjecture, cleaned question and verification metadata.
\item \emph{Geometric Group Theory, Hyperbolic Dynamics and Symplectic Geometry (2007)}, Oberwolfach Reports, pp.\ 1991--2058, DOI 10.4171/owr/2006/33: \url{https://doi.org/10.4171/owr/2006/33}. The citation uses 2007, while the DOI and proposed-year metadata use 2006.
\item Supplied reference ``Lifschitz--Witte Morris'': \url{https://arxiv.org/abs/math/0604612}.
\item Literature triage checked 2026-08-21; local record updated 2026-09-16. No external source has been visited or independently checked here.
\end{itemize}
\end{document}
